\documentclass[11pt]{article}
\usepackage[margin=2.3cm]{geometry}
\usepackage{amsmath,amssymb,booktabs}
\usepackage[colorlinks,linkcolor=blue,citecolor=blue]{hyperref}
\newcommand{\so}{\mathrm{so}(3)_k}
\newcommand{\Z}{\mathbb{Z}}
\newcommand{\Q}{\mathbb{Q}}
\title{Selection rules, entropic order parameters and the massive sector\\ of the $N=2$ flow $k\to k-2$}
\date{Working note}
\begin{document}
\maketitle

\noindent\textbf{Scope.} Bosonic $N=2$ coset $(l,m,s)$ with the conventions of \texttt{N2minimal.m} (S-matrix eq.~(2.6) of N2Gab normalised by $1/K$, $K\equiv k+2$). I take $k$ odd throughout (no subtleties with the invertible line $[k,0,0]$). For a node $T$ I compute the commuting lines $F_T$ and $\mu_T=(\sum_{F_T}d^2)^2$ (paul.pdf eq.~(5)); $\lambda_T\equiv\sqrt{\mu_T}=\dim F_T$ is the index of $T$ in the completion and the constant in the two-interval certainty relation $S_{\rm int}+S_{\rm tw}=\log\lambda_T$. The numbers come from \texttt{n2\_nodes.py}, which I checked for $k=3,\dots,11$.

\section{UV node of $(0,-2,0)$ and IR node of its image}

\paragraph{UV.} The smallest node containing $\phi_{UV}=(0,-2,0)\sim(k,k,2)$ (and its conjugate) is
\begin{equation}
N_{UV}=\langle(0,-2,0)\rangle=\{(0,2n,0)\}_{n\in\Z_K}\simeq \Z_K ,\qquad F_{UV}=\{(L,0,S)\}_{L+S\ {\rm even}}\simeq \so\boxtimes\Z_4 ,
\end{equation}
with $\Z_4=\{(0,0,0),(0,K,\pm1),(0,0,2)\}$ generated by $g=(0,K,1)$, $g^2=(0,0,2)$. Hence
\begin{equation}
\mu_{UV}=16\,(\dim\so)^2,\qquad \dim\so=\sum_{L\ {\rm even}}d_L^2=\frac{K}{4\sin^2(\pi/K)} .
\end{equation}

\paragraph{IR.} The IR is approached by the D-term of the non-BPS field $(4,\pm2,0)$ of the $k'=k-2$ model (N2Gab \S5.3), with $\Delta=2+10/k$. This field exists for $k'\ge4$. Its node is
\begin{equation}
N_{IR}=\langle(4,2,0)\rangle_{k-2}=E_{k-2},\qquad E_{k}\equiv\{(l,m,0):\ l,m\ \text{even}\},\qquad F_{IR}=\Z_4,\qquad \mu_{IR}=16 .
\end{equation}
The only $k-2$ nodes that contain $(4,2,0)$ have $\mu\in\{1,4,16\}$.

\begin{center}
\begin{tabular}{ccccccc}
\toprule
$k$ & $\dim\so$ & $\mu(N_{UV})$ & $\mu(E_k)$ & $\mu(N_{IR})$ & $h_g$ (UV) & $h_{g'}$ (IR) \\
\midrule
3 & 3.618 & 209.44 & 16 & 16$^*$ & 7/8 & 3/8\\
5 & 9.296 & 1382.6 & 16 & 16$^*$ & 3/8 & 7/8\\
7 & 19.23 & 5919.4 & 16 & 16 & 7/8 & 3/8\\
9 & 34.65 & 19206 & 16 & 16 & 3/8 & 7/8\\
11 & 56.75 & 51523 & 16 & 16 & 7/8 & 3/8\\
\bottomrule
\end{tabular}\\[2pt]
{\small $^*$For $k\le5$, $(4,2,0)$ does not exist; the entry is the node $E_{k-2}$.}
\end{center}

\noindent $\mu(N_{UV})$ matches no node of the $k-2$ model, as you found. The two ingredients of the mismatch behave differently. $\so$ lives in $\Q(\zeta_K)$ and the IR in $\Q(\zeta_{8(K-2)})$, and these fields intersect only in $\Q$. $\Z_4$ is the only piece that can be matched.

\section{No-go: $\so$ cannot act on an IR that contains the $k-2$ model}

\textbf{Claim.} For odd $k\ge3$, no IR phase of the form $\mathrm{CFT}_{k-2}\oplus(\text{gapped vacua})$ carries the lines $\so$, with any number of gapped vacua.

\emph{Argument.}
\begin{enumerate}
\item No topological interface connects $c_{k-2}>0$ with $c=0$. A line therefore maps the CFT vacuum to itself, $L|\Omega_{\rm CFT}\rangle=\chi(L)|\Omega_{\rm CFT}\rangle$. Its restriction to the CFT block is nonzero, since $L\bar L\ni1$, and is a sum of lines of the $k-2$ model. So $\chi(L)=\sum_j n_j d_j>0$.
\item $\chi$ is then a character of the $\so$ fusion ring that is positive on every basis element. Such a character must be $\mathrm{FPdim}$. Hence $d_L\in\Q(\zeta_K)\cap\Q(\zeta_{8(K-2)}\,{\rm or}\,\zeta_{24\cdot 8(K-2)})$, whose real part is $\Q$ because $\gcd(K,K-2)=1$ with $K$ odd.
\item But $d_2=1+2\cos(2\pi/K)$ is irrational for odd $K\ge5$. Contradiction.
\end{enumerate}
The Witten index fixes the number of CFT blocks: $k+1=(k-1)+2$. A tensor-product IR $\mathrm{CFT}\otimes\mathrm{TQFT}_n$ would give $n(k-1)$ Ramond ground states and is excluded.

\textbf{Consequence (selection.pdf, transportability).} Any flow that stays inside $N_{UV}$ preserves $F_{UV}$ exactly, since the lines commute with every operator of the node, the perturbed stress tensor included. Such a flow is therefore fully massive. Section 4 shows that this is the Chebyshev flow.

\section{The massless flow lives in a larger node: $E_k$, with $\mu=16$}

The LG coupling of $X^{k}$ has dimension $y_k=2/K$, and the F-term of $X^{k-2}$, i.e.\ $(k-2,k-2,2)\sim(2,-4,0)$, has $y_{k-2}=4/K=2y_k$. These are resonant, so at $O(\lambda^2)$ the linear superpotential $X^{k+2}+\lambda X^k$ differs from the pure $\phi_{UV}$ perturbation by a coupling $g\propto\lambda^2$ to $(2,\mp4,0)$. That operator is charged under $\so$, since $(L,0,0)$ acts on $l=2$ by a ratio $\neq d_L$. This is the CFT version of the second-order obstruction in N2Gab (\S3.4, \S5.2).

The Chebyshev superpotential is the unique completion that keeps $\so$ (N2Gab \S4), and $\lambda$ is its flat coordinate. The flow inside $N_{UV}$ is therefore the Chebyshev one. The massless flow is instead generated by
\begin{equation}
N'_{UV}=\langle(0,-2,0),(2,-4,0)\rangle=E_k,\qquad F'=\Z_4,\qquad \mu'=16=\mu(E_{k-2})=\mu(N_{IR}).
\end{equation}
So the selection rule is satisfied as $E_k\to E_{k-2}$, and the chain $N_{UV}\subset E_k$ has relative index $\lambda=\dim\so$.

\emph{Anomaly matching.} The Verlinde lines of the simple current $g=(0,K,1)$ form $\mathrm{Vec}^{\omega}_{\Z_4}$. For $K$ odd, $h_g=(1-2K)/8$ mod 1, which is $7/8$ or $3/8$. So $\omega=8h_g$ mod 4 equals $3$ for every odd $k$, in both the UV and the IR (table). The generator of $\Z_4$ is anomalous, while $\Z_2=\langle(0,0,2)\rangle$, with $h=1/2$, is anomaly free. The DHR categories of the two nodes, $\mathcal Z(\mathrm{Vec}^\omega_{\Z_4})$, therefore coincide.

On the chiral ring, $g$ acts as $(-1)^l$, i.e.\ $X\to-X$. Together with the $R$-rotation it squares to $(-1)^{s}$. For odd $d=K$ this is the LG symmetry $X\to -X$, $W\to -W$, compensated by $\theta\to i\theta$. It preserves $X^d+\lambda X^{d-2}$ and exchanges the massive vacua $X=\pm i\sqrt\lambda$.

\section{TQFT content from entropic order parameters}

Work in the node $E_k$, with $\lambda=4$, and use the certainty relation of entropic.pdf (3.55)/(3.60) together with its scaling analysis (\S3.4, \S4):
\begin{equation}
S_{\rm int}(x;R)+S_{\rm tw}(x;R)=\log 4 .
\end{equation}

\begin{itemize}
\item \emph{CFT vacuum.} All sectors of $E_k$ stay massless, because $\mu(E_k)=\mu(E_{k-2})$. The massless subalgebra of massless.pdf contains all of $E_{k-2}$. Both parameters approach conformal functions of $x$, with $U(1)=-\tfrac12\log 16$ (modular.pdf (3.30)). No TQFT is required on this block.

\item \emph{Gapped vacua.} Every sector of $E_k$ is massive. By \S4 of entropic.pdf, each order parameter tends to a shape-independent constant, $S_{\rm int}\to\log\lambda_{\rm SSB}$ and $S_{\rm tw}\to\log\lambda_{\rm unbr}$, with
\begin{equation}
\lambda_{\rm SSB}\,\lambda_{\rm unbr}=4 .
\end{equation}
An unbroken part must be anomaly free (it needs a fiber functor). Since $\omega(\Z_4)\neq0$ and $\omega|_{\Z_2}=0$, the options are $\lambda_{\rm SSB}\in\{2,4\}$. This gives the bound on the TQFT total dimension
\begin{equation}
2\le D^2_{\rm TQFT}=\lambda_{\rm SSB}\le 4=\sqrt{\mu(E_k)} .
\end{equation}
The upper bound is the certainty-relation bound $S_{\rm int}\le\log\lambda$. The lower bound is anomaly inflow.

\item The LG count of massive vacua, 2, saturates the lower bound. \textbf{The TQFT is the $\Z_4\to\Z_2$ SSB phase.} It has two vacua exchanged by $g$, unbroken $(-1)^F=(0,0,2)$, and topological local operators $\{1,\sigma\}$ with $\sigma\sim\mathrm{sign}(X)$, $d_\sigma=1$, $D^2=2$. In these vacua $S_{\rm int}\to\log2$ (constant law) and $S_{\rm tw}\to\log2$.
\end{itemize}

\paragraph{Check with the Chebyshev (pure-$N_{UV}$) flow.} Here $\lambda=4\dim\so$. $\so$ has no fiber functor, since its dimensions are not integers, so it must be broken. The number of vacua $n$ must satisfy $n\ge\deg\Q(d_2)=\varphi(K)/2$, because a rank-$n$ NIM-rep needs eigenvalues of degree $\le n$. $\Z_4$ is broken to $\Z_2$ as above. The regular module then gives $n=\tfrac{k+1}{2}\times2=k+1$ vacua and $\lambda_{\rm SSB}\lambda_{\rm unbr}=\dim\so\cdot2\cdot2$. This reproduces exactly the $k+1$ massive vacua of the Chebyshev flow (N2Gab \S4, the Witten index). The same analysis also shows why only 2 gapped vacua cannot host $\so$ for $k\ge5$: $\varphi(K)/2>2$.

\section{Relation to massless.pdf}

In each vacuum sector, take $W_m$ = operators with $\mu_O\ge m>0$ and $A_m$ its commutant-type algebra.
\begin{itemize}
\item In the CFT vacuum of the massless flow, no sector of $E_k$ is in $W_m$, so $A_m\supseteq E_{k-2}$ and no emergent symmetry is needed. The gap condition of the notes, no accumulation of masses near 0, holds: the only masses are $\sim|\lambda|^{K/4}$ from kinks to the massive block.
\item The conjecture ``massless and massive together $\Rightarrow$ symmetry'' is realised here across vacua, not within one. The symmetry is the anomalous $\Z_4$. Its anomaly forbids the massive block from being a single trivially gapped vacuum and forces the two-vacuum TQFT.
\item For the pure $N_{UV}$ flow, the argument of \S2 shows that no $\so$-charged sector can remain massless. If one did, the neutral stress tensor would leave a CFT block on which $\so$ would have to act. So $A_m$ is trivial and the flow is gapped.
\end{itemize}

\section*{Summary}
\begin{enumerate}
\item $\langle(0,-2,0)\rangle=\Z_K$ with $F=\so\boxtimes\Z_4$ and $\mu=16\dim^2\so$. No $k-2$ node matches it, and \S2 shows that no IR containing the $k-2$ model can carry $\so$. Staying in this node means the Chebyshev massive flow, with $k+1$ vacua $=$ regular $\so$ module $\times$ $\Z_4/\Z_2$.
\item The linear (massless) LG flow resonantly switches on $(2,\mp4,0)$ at $O(\lambda^2)$. The node becomes $E_k=\{(\text{even},\text{even},0)\}$ with $F=\Z_4$ and $\mu=16$. It flows to $E_{k-2}=\langle(4,2,0)\rangle$, with $\mu=16$ and matching $\Z_4$ anomaly $\omega=3$.
\item Entropic bound on the TQFT: $2\le D^2_{\rm TQFT}\le4$. The LG count gives $D^2=2$: the two massive vacua are the $\Z_4\to\Z_2$ SSB TQFT of the anomalous $\Z_4$ generated by $(0,K,1)$ ($X\to-X$ with $R$-rotation).
\end{enumerate}

\paragraph{Caveats.}
\begin{itemize}
\item Odd $k$ only. For even $k$, $[k,0,0]$ is invertible and survives (N2Gab App.~D), so $F'$ is larger.
\item The identification ``flow inside $N_{UV}$ $=$ Chebyshev'' uses the uniqueness of N2Gab \S4 within their ansatz, plus flat-coordinate $=$ CFT-coupling.
\item The field-theoretic statement in \S2 assumes that the IR lines commuting with the image node have dimensions in the cyclotomic field of the $k-2$ model, which holds for lines of that rational model.
\item The anomaly class is read off from the simple-current spin, $\omega=8h_g$ mod 4.
\end{itemize}

\end{document}
